\section{近期研究（2026 年 8--11 月）}
\paragraph{XClipGS：半空间精确积分与切面监督}
2026.08~\cite{2608_07760}。交互式切割需要处理横跨平面的高斯。按中心点硬裁会把整个原语保留或删除，容易在切面附近留下错误贡献。XClipGS 在仿射 EWA 模型下直接积分高斯落在保留半空间的部分，并用同一个可微算子训练。切割视图覆盖不同平面位置和相机距离，为原本被外层遮住的体内结构提供监督。

作者在六个 CT 和两个 MRI 体数据上比较切面边缘和错误侧漏。实验包括在同一个体内表示上替换裁剪算子，以区分算子与训练监督的作用。算子可计算新的半空间，无需再次训练，但切面视觉质量仍依赖体内表示。留出偏移位于训练区间内，这组结果主要支持对该区间内新切面的泛化。

\paragraph{用多色 X 射线投影降低锥束 CT 金属伪影}
2026.08~\cite{2608_13159}。锥束 CT 的射线跨越金属时，低能成分往往比高能成分衰减更多，单能模型难以解释随路径变化的衰减。该方法用高斯基元表示局部体结构，以每个高斯的两个系数分别表达近似与能量无关的康普顿项和随能量变化的光电项。可微多色投影将不同能量下的透射信号按系统响应合成，并从实际投影联合估计体场和全局系统响应。

原文将肺、牙和西兰花三个模拟场景分别嵌入铁、钛、铝物体，按体 PSNR 与 SSIM 对照重建，并对五组实际扫描作比较。实际样本的对照条件各异，例如金属杆样本未取得单独的无金属参考。结果支持所测扫描条件下的伪影抑制，不能直接外推到其他扫描设备、材料组合或临床应用。

\paragraph{TR-GS：用重尾原语和射线置信度重建稀疏视角 CT}
2026.08~\cite{2608_16042}。TR-GS 用可投影的 Student t 分布取代放射高斯，以缓解稀疏观测和异常值对体密度估计的影响。局部射线的几何支撑决定原语自由度的正则目标。置信度同时引导三维小波约束，在压制稀疏采样伪影时保留可辨识的高频边界。

作者使用 R²-Gaussian 的合成和真实 CT 数据，在整圈扫描轨迹上均匀选取 12、18 和 25 个视角比较，另外用 6 和 50 个视角做消融。论文报告大部分设置下改善最显著的是少视角情形。整圈均匀稀疏与只覆盖窄角度属于不同的逆问题，后者留待验证。

\paragraph{统计形状约束下的少视角 X 射线股骨重建}
2026.09~\cite{2609_10376}。少视角 X 射线中，不同的三维密度可能产生近似投影。该方法先用群体平均模板初始化高斯的位置、取向和密度，再把原语的位置绑定到统计形状模型，让它们按群体上可见的形变模式共同移动。密度先验另负责约束体强度。

论文从六例留出股骨 CT 合成投影，在五视角下相对 R²-Gaussian 的平均 PSNR 从 34.55 升至 37.38 dB，十视角下提升较小。先验在投影约束最少时作用更明显，结果目前限于该股骨模板和模拟投影设置。

\paragraph{LINGO：改善稀疏 X 射线重建的初始化和低密度梯度}
2026.09~\cite{2609_22849}。LINGO 在 R2-Gaussian 上从投影视图的掩码构造三维体素筛选，抑制背景噪声并初始化高斯。训练用动态损失缩放加强低密度区域的梯度，并以 IPSD 专门评价初始点云结构。

论文在 X3D 比较 5、10 和 25 张 X 射线输入。结果对应数据集的器官和投影条件；IPSD 衡量初始化形态，不能代替最终 CT 体素重建或临床评价。

\paragraph{FactorSplat：保持体渲染代理几何并编辑传递函数}
2026.10~\cite{2610_02382}。体渲染的区域传递函数把组织强度映射为颜色和不透明度。FactorSplat 保留一套高斯位置、协方差和方向相关响应，把传递函数的变化作用于视角无关颜色和不透明度。局部组织与强度查表给出物理外观改动，低秩残差弥合代理表示和原体渲染器间的差异。修改传递函数后，交互转动相机可复用更新后的高斯属性。

作者在七个 CT 与 MRI 扫描中比较留出的传递函数，报告相对于区域感知 VEG 的平均重建和编辑精度提升。逐扫描结果仍有例外，例如心脏及膝关节的编辑误差由 VEG 取得更低值。模型对所支持的传递函数族训练，更广泛的编辑需要单独验证。

\paragraph{CellSplat4D：将显微成像轴向模糊与稀疏时间补全分开}
2026.10~\cite{2610_04199}。CellSplat4D 将遗漏的活细胞显微体数据视为时空补全。它用物理测量固定轴向点扩散函数，避免优化误把光学模糊归咎于原语伸展；时间门以帧为单位，在固定预算下按门控残差增密。连续时间表示可用来渲染未采样时刻。

CE02 和 CHO 的比较让各方法共享同一前向光学模型，并分别检查细胞核、留出轴向层等指标。作者发现离最近采集帧越远时本法的保真度衰减较慢，但部分中心定位指标仍落后。结果说明特定采集协议下的补全表现，未证明任意动态都能忠实恢复。

\paragraph{AD-GS：把稀疏牙科 CBCT 的原语分配与结构约束分开}
2026.10~\cite{2610_14963}。AD-GS 从 FDK 初值建立放射高斯，用衰减密度引导马尔可夫链蒙特卡罗采样，将更多原语分配给牙齿及颌面硬组织。基于 MedSAM 的语义约束负责维持结构一致性，分阶段训练把重点从整体形状转向局部细节。投影一致性仍是重建的主要观测约束。

论文在 ToothFairy3 上用 TIGRE 模拟含噪的 20、30、50 个圆轨迹投影。实验同时报告体图像质量、分割重叠和表面距离，因而能观察细边界是否随重建改善。该协议支持评价模拟少视角条件，真实低剂量扫描仍需相应验证。

\paragraph{面向固定高斯预算的跨个体解剖分配}
2026.10~\cite{2610_20397}。对已知 CT 体进行可交互渲染时，固定预算里每个高斯放在哪里会影响解剖边界的表现。作者从训练个体逐例优化的高斯模型提取每个体素格的典型分配，训练标签感知的神经分配器，并给出一个考虑可见性的解析方案。部署时先按分配密度抽取恰好 B 个中心，再对当前体数据拟合颜色、尺度等属性。

作者比较相同预算下的均匀分配、剪枝和所提方法，在 20 例留出和 19 例外部个体上评估。评价分配器时固定高斯中心，以隔离位置的作用。两种分配器各有外观适用范围：解析法依赖解剖标签仍能解释渲染外观，学习法也受训练时传递函数对比度的影响。

\paragraph{KFS⁴LAM：让内镜关键帧筛选与形变建图协同}
2026.10~\cite{2610_22505}。内镜视频既含冗余相邻帧，也含局部组织运动。KFS⁴LAM 按累计运动提出关键帧候选，再根据图像冗余筛选。建图端按保留帧调节优化预算，用解剖褶皱边界监督与受限时间形变维持结构一致性。

作者在 C3VD 和 C3VD-v2 上以轨迹误差、图像相似度和深度误差评估，并在 SplaTAM 与 EndoGSLAM 主干下做匹配比较。自适应优化预算会增加部分条件下的计算量。作者将结果定位为质量与计算的取舍，并指出外部褶皱检测器的可靠性会影响推广。

\paragraph{X-Splat：用同一组高斯重建 RGB 外观和 X 射线体密度}
2026.10~\cite{2610_23237}。X-Splat 面向同一物体的标定 RGB 图像和 X 射线投影，先用粗 X 射线重建建立跨模态配准和共享几何初值。优化时两种模态共享一组高斯几何，颜色等外观属性集中在近表面原语，X 射线重建给出用于 RGB 的表面深度约束。

实验包括 InsideOut 的四个场景和 Flower \& Insect 的五个场景，并从参考 CT 合成统一协议下的锥束投影。冻结 CT 并单独优化 RGB 的消融不及联合模型。作者报告联合训练主要改善 RGB，X 射线质量约等于单模态基线。透明物体或强次表面散射会破坏近表面外观假设。

\paragraph{H-2DGS：用高度场高斯表达皮层表面}
2026.10~\cite{2610_30910}。H-2DGS 先建立体高斯，再从视图深度融合出截断符号距离场，最后在零水平集上布置曲面高度场原语。高度场拟合局部皮层曲率，颜色和不透明度基函数保留纹理变化。术后配准以 MRI 网格渲染的深度约束目标形状，同时使用颜色和重叠约束维持表面外观。

论文在五例中报告更低的多视图深度不一致，但 PSNR 与 SSIM 逊于体高斯基线。另一个视频中的生理脉动影响 SfM 位姿，导致后续重建失败。结果体现了显式薄表面带来的几何优势及其对外观表达与输入位姿的要求。

\paragraph{CTBakeGS：从已知 CT 体直接烘焙渲染代理}
2026.10~\cite{2610_37406}。CTBakeGS 处理的是已有 CT 密度体及选定传递函数的显示。它直接将体数据和渲染外观转换为标准三维高斯，使用毫米单位的自适应八叉树划分，并在转换中补偿透明度和抗锯齿。输出可用普通高斯查看器显示，转换时无需先渲染训练相机或逐例拟合。

作者在六个 TotalSegmentator v2 扫描、三种传递函数预设下与参考体渲染比较，报告宽视图均方根误差约为满量程的 2.3\% 至 3.1\%。这衡量的是对既有渲染外观的复现，适用范围与传递函数和参考渲染配置绑定。

\paragraph{CryoGaussian：用显式密度高斯表示冷冻电镜异质性}
2026.11~\cite{2611_26128}。异质冷冻电镜重建需要同时估计粒子状态和每种状态的三维密度。CryoGaussian 从随机初始化的原语出发，用粒子与高斯层级的潜变量调节中心、幅值和尺度，在可微加和投影中拟合粒子图像。这里高斯用于表达密度，优化利用的是冷冻电镜观测模型。

作者按 CryoBench 的逐构象和逐图像协议比较傅里叶壳相关，并在真实粒子数据上检验重建。方法输入包含固定的粒子朝向、平移和对比传递函数。真实实验数据没有真值，论文采用已有参考重建等证据评价，证据条件不同于模拟基准。
